\documentclass[12pt]{article}

\usepackage{sbc-template}
\usepackage{graphicx,url}
\usepackage[utf8]{inputenc}
\usepackage[brazilian]{babel}
\usepackage{minted}
\usepackage{float}
\usepackage{subcaption}
\usepackage{booktabs}
\usepackage{url}
\usepackage{pdflscape}
\usepackage{hyperref}
\usepackage{tikz}
\usetikzlibrary{positioning, arrows.meta, shapes.geometric}

\renewcommand{\listingscaption}{Listagem}

\NewCommandCopy{\svtt}{\texttt}
\renewcommand\texttt[1]{\svtt{\small#1}}

\setminted{
  style=friendly,        % tema de cores
  fontsize=\footnotesize,% tamanho adequado para colunas SBC
  % linenos=true,          % numeração de linhas
  breaklines=true,       % quebra linhas longas
  frame=lines,           % linhas acima e abaixo
  framesep=4mm
  % baselinestretch=1.1,
  % xleftmargin=2em,
  % numbersep=6pt,
  % numberfirstline=true
}
     
\sloppy

\title{Jass (Jaime's home ASsistant Script)}

\author{Jaime Antonio Daniel Filho\inst{1} }

\address{Curso de Ciência da Computação -- \\ Universidade Federal de Santa Maria (UFSM)
  \email{jafilho@inf.ufsm.br}
}

\begin{document} 

\maketitle

\section{Introdução}

O Home Assistant consolidou-se como uma das principais plataformas de automação residencial, permitindo integrar e controlar dispositivos de diferentes fabricantes em um único ambiente. Entretanto, a criação de automações é baseada em arquivos YAML (\emph{YAML Ain't Markup Language}), cuja estrutura pode se tornar extensa e complexa mesmo para tarefas simples. Uma automação básica, como acender uma luz ao abrir uma porta, exige o conhecimento de diversos elementos da sintaxe, como gatilhos, condições, ações, serviços e parâmetros específicos de cada dispositivo. Além disso, muitos erros só são percebidos durante a execução da automação, seja por um nome de serviço incorreto, um parâmetro incompatível ou um tipo de valor inválido, dificultando a identificação da causa do problema e tornando o processo de desenvolvimento menos acessível para usuários sem experiência técnica.

Com o objetivo de reduzir essa complexidade, este trabalho apresenta o Jass, uma \emph{Domain-Specific Language} (DSL) capaz de transpilar descrições de alto nível para o formato YAML utilizado pelo Home Assistant. A proposta é oferecer uma sintaxe mais simples e intuitiva, abstraindo detalhes estruturais da plataforma e permitindo que o usuário descreva apenas o comportamento desejado da automação. A linguagem também incorpora mecanismos de validação léxica, sintática e semântica, possibilitando que erros, como construções inválidas ou operações incompatíveis com determinados dispositivos.


A Listagem~\ref{lst:jass_example1} e a Listagem~\ref{lst:jass_example1_yaml} ilustram, respectivamente, a automação completa em
Jass e o YAML equivalente que o compilador gera. Note como \texttt{turn\_on(light.varanda,\ ...)} se
expande para o serviço \texttt{light.turn\_on} com a estrutura
\texttt{target}/\texttt{data}, e como \texttt{delay(5m)} vira a string
\texttt{"00:05:00"} no formato exigido pela plataforma.

\begin{listing}[H]
    \centering
    \begin{minted}{swift}
automation "Luz da varanda" {
    when (binary_sensor.porta changes to "on")
    do {
        turn_on(light.varanda, brightness=180)
        delay(5m)
        turn_off(light.varanda)
    }
}
    \end{minted}
    
    \caption{Uma automação básica em Jass.}
    \label{lst:jass_example1}
\end{listing}

\begin{listing}[H]
    \centering
    \begin{minted}{yaml}
- id: "1780962966"
  alias: "Luz da varanda"
  mode: "single"
  triggers:
    - trigger: state
      entity_id: binary_sensor.porta
      to: "on"
  actions:
    - action: light.turn_on
      target:
        entity_id: light.varanda
      data:
        brightness: 180
    - delay: "00:05:00"
    - action: light.turn_off
      target:
        entity_id: light.varanda
    \end{minted}
    
    \caption{A saída YAML da automação anterior.}
    \label{lst:jass_example1_yaml}
\end{listing}

O restante do relatório está organizado da seguinte forma: a Seção 2 apresenta as instruções para execução; a Seção
3 descreve a arquitetura; a Seção 4 define a
linguagem de forma conceitual; a Seção 5 a formaliza como gramática
livre de contexto; as Seções 6 a 8 detalham as fases de análise; a Seção 9 trata da geração de código; a Seção 10 reúne
exemplos completos e a Seção 11 conclui, sintetizando os resultados.

\section{Instruções para execução}

O projeto usa o gerenciador \href{https://docs.astral.sh/uv/getting-started/installation/}{uv}, que pode ser instalado em Linux, macOS e Windows. A partir do diretório do
trabalho, instale as dependências e o comando Jass (Listagem~\ref{lst:uv_install}):

\begin{listing}[H]
    \centering
    \begin{minted}{bash}
uv sync
\end{minted}
    \caption{Comando para configurar o projeto.}
    \label{lst:uv_install}
\end{listing}

Para transpilar um arquivo \texttt{.jass} para YAML (Listagem~\ref{lst:uv_run}):

\begin{listing}[H]
    \centering
    \begin{minted}{bash}
uv run jass caminho/para/automacao.jass
\end{minted}
    \caption{Comando para transpilar um arquivo .jass para YAML.}
    \label{lst:uv_run}
\end{listing}

Para gravar o resultado em arquivo, ao invés de mostrar no terminal, com \texttt{-o} (Listagem~\ref{lst:uv_output}):

\begin{listing}[H]
    \centering
    \begin{minted}{bash}
uv run jass caminho/para/automacao.jass -o automacao.yaml
    \end{minted}
    \caption{Comando para gravar o resultado em arquivo.}
    \label{lst:uv_output}
\end{listing}

Para inspecionar uma fase intermediária do compilador, com
\texttt{-\/-emit} (Listagem~\ref{lst:uv_emit}), especificando \texttt{tokens}, \texttt{parse},
\texttt{ast}, \texttt{check} ou \texttt{yaml}:

\begin{listing}[H]
    \centering
    \begin{minted}{bash}
# imprime a AST tipada
uv run jass caminho/para/automacao.jass --emit ast   

# roda a análise semântica
uv run jass caminho/para/automacao.jass --emit check 
    \end{minted}
    \caption{Comando para inspecionar uma fase intermediária do compilador.}
    \label{lst:uv_emit}
\end{listing}

Em caso de erro, os diagnósticos são impressos no \texttt{stderr} e o
programa termina com código não nulo.

\section{Visão geral da arquitetura}

O Jass é organizado como um compilador em \emph{pipeline}, ou
seja, em uma sequência de fases em que cada uma consome o artefato
produzido pela anterior e gera o artefato da seguinte. Essa estrutura
mantém as responsabilidades separadas e torna o sistema fácil de testar
e inspecionar fase a fase. A Tabela~\ref{tab:fases_modulos}
associa cada fase ao módulo que a implementa e ao artefato que ela
produz.

% \begin{figure}[htpb]
%     \centering
%     \begin{tikzpicture}[
        
%         node distance = 1.2cm,
%         auto,
        
%         box/.style = {
%             rectangle, 
%             draw=black, 
%             thick,
%             align=center, 
%             minimum width=4.5cm, 
%             minimum height=0.8cm, 
%             rounded corners=3pt
%         },
        
%         arrow/.style = {
%             -{Stealth[scale=1.2]}, 
%             thick
%         }
%     ]

    
%     \node [box] (src) {Texto-fonte\\(.jass)};
%     \node [box, below=of src] (tok) {Tokens};
%     \node [box, below=of tok] (tree) {Árvore de\\derivação};
%     \node [box, below=of tree] (ast) {AST tipada};
%     \node [box, below=of ast] (chk) {Diagnósticos +\\tabela de símbolos};
%     \node [box, below=of chk] (model) {Modelo de saída\\(dicts/listas)};
%     \node [box, below=of model] (yaml) {YAML do\\Home Assistant};

    
%     \draw [arrow] (src) -- node[right] {análise léxica} (tok);
%     \draw [arrow] (tok) -- node[right] {análise sintática} (tree);
%     \draw [arrow] (tree) -- node[right] {transformação} (ast);
%     \draw [arrow] (ast) -- node[right] {análise semântica} (chk);
%     \draw [arrow] (chk) -- node[right] {geração de código} (model);
%     \draw [arrow] (model) -- node[right] {serialização} (yaml);

%     \end{tikzpicture}
    
    
%     \caption{Fluxo do compilador.}
%     \label{fig:fluxo_compilacao}
% \end{figure}

\begin{table}[H]
\centering
\begin{tabular}{l l l}
\toprule
\textbf{Fase} & \textbf{Módulo}\\
\midrule
Análise léxica & \texttt{parser.py} (Lark) \\
Análise sintática & \texttt{parser.py} (Lark, LALR) \\
Transformação & \texttt{transform.py}, \texttt{ast.py}, \texttt{values.py} \\
Análise semântica & \texttt{analyzer.py}, \texttt{registry.py} \\
Geração de código & \texttt{codegen.py} \\
Orquestração / diagnósticos & \texttt{cli.py}, \texttt{errors.py}  \\
\bottomrule
\end{tabular}
\caption{Fases do compilador e módulos correspondentes.}
\label{tab:fases_modulos}
\end{table}


Os artefatos intermediários são, em ordem: o texto-fonte; a sequência de tokens, unidades léxicas como
palavras-chave, identificadores e literais; a árvore de derivação
(\emph{parse tree}), que reflete a estrutura gramatical da
entrada; a AST (\emph{Abstract Syntax Tree}), uma representação
tipada e enxuta, na qual cada nó é um objeto Python com tipo próprio e
posição de origem; o modelo de saída, uma árvore de dicionários e listas
Python já na forma esperada pelo Home Assistant; e, por fim, o YAML
serializado. A análise semântica é uma fase à parte: percorre a AST para
preencher a tabela de símbolos e validar a entrada, produzindo a decisão de que o programa está correto ou a
lista de erros.

% \subsection{A ferramenta Lark}

As duas primeiras fases são construídas sobre o Lark, uma biblioteca
Python para construção de analisadores. O Lark recebe uma gramática em
notação EBNF e oferece, a partir dela, o \emph{lexer} (analisador
léxico) e o \emph{parser} (analisador sintático), além de utilitários
para o tratamento de erros. O Jass o emprega na configuração
LALR(1) com \emph{lexer} contextual, aproveitando:

\begin{itemize}
\tightlist
\item
  a propagação automática de linha e coluna para cada nó, base dos
  diagnósticos localizados;
\item
  as exceções de erro distintas para caractere inválido e \emph{token}
  inesperado, que sustentam a separação entre erro léxico e sintático;
\item
  o \emph{parser} interativo, que permite alimentar \emph{tokens}
  manualmente e, assim, implementar a recuperação de erros em modo
  pânico.
\end{itemize}

A gramática reconhece apenas a forma das construções, isto é, a ordem das
seções, o formato de um gatilho, a estrutura de uma chamada, enquanto
regras que dependem de tipos, domínios e nomes ficam a cargo do
analisador semântico, que tem mais contexto para produzir mensagens
precisas.

\section{Definição da Linguagem}

O projeto da linguagem foi orientado por dois princípios fundamentais. O primeiro é fornecer uma sintaxe declarativa e uniforme, estruturada em seções com responsabilidades bem definidas. O segundo é adotar um vocabulário reduzido e composto por nomes amigáveis, como \texttt{turn\_on}, \texttt{delay} e \texttt{say}, cujo significado é conhecido e validado pelo compilador. Esta seção apresenta, em nível conceitual, as principais construções oferecidas pela linguagem.

\subsection{Automação}

A unidade básica de um programa é a automação, introduzida pela
palavra-chave \texttt{automation} seguida de um título e de um corpo
entre chaves. Um arquivo contém zero ou mais automações e,
opcionalmente, definições de nomes. O corpo é composto
por seções em ordem fixa: metadados opcionais
(\texttt{id}, \texttt{description}, \texttt{mode}, \texttt{max}), a seção
\texttt{when} (gatilhos), a seção \texttt{if} (condições) e a seção \texttt{do} (ações). A
Listagem~\ref{lst:skeleton} apresenta o esqueleto de uma automação com
as quatro seções. A ordem metadados, \texttt{when}, \texttt{if} e
\texttt{do} é imposta pela gramática.

\begin{listing}[H]
    \centering
    \begin{minted}{swift}
automation "Aquecedor da suíte" {
    mode = "single"                         // metadado
    when (switch.toalheiro changes to "on") // gatilho
    if (between @15:00:00 @23:00:00)        // condição
    do { turn_on(switch.toalheiro) }        // ações
}
    \end{minted}
    \caption{Esqueleto de uma automação.}
    \label{lst:skeleton}
\end{listing}

\subsection{Gatilhos}

A seção \texttt{when} lista um ou mais gatilhos, os eventos
que disparam a automação, separados pela palavra \texttt{or}. Cada
gatilho pode receber um identificador com \texttt{as\ "id"},
referenciado depois nas condições. Há quatro famílias, exemplificadas na Listagem~\ref{lst:triggers}:

\begin{itemize}
\item \textbf{de estado:} \texttt{ENTIDADE changes [from V] [to V] [for DURAÇÃO]}
  reage à mudança de estado de uma entidade, opcionalmente filtrando
  origem/destino e exigindo permanência por um tempo;
\item \textbf{numérico:} \texttt{ENTIDADE (\textless{} \textbar{} \textgreater{}) NÚMERO [for DURAÇÃO]}
  reage ao cruzamento de um limiar;
\item \textbf{de horário:} \texttt{@HH:MM:SS} dispara em um horário do
  dia;
\item \textbf{solar:} \texttt{(sunset \textbar{} sunrise) [(+ \textbar{} -) DURAÇÃO]}
  dispara no nascer ou pôr do sol, com deslocamento opcional.
\end{itemize}

\begin{listing}[H]
    \centering
    \begin{minted}{swift}
when (
    binary_sensor.sala changes to "on" for 5m as "movimento" or
    sensor.temperatura > 30 as "calor" or
    @07:00:00 as "manha" or
    sunset - 30m as "anoitecer"
)
    \end{minted}
    \caption{Exemplos de gatilhos.}
    \label{lst:triggers}
\end{listing}

\subsection{Condições}

A seção \texttt{if} é uma expressão booleana que precisa ser
verdadeira para a automação prosseguir. As condições abrangem:

\begin{itemize}
  \item \textbf{estado:} \texttt{ENTIDADE == "v"} ou \texttt{!=};
  \item \textbf{numérica:} \texttt{ENTIDADE \textless{} N} ou \texttt{\textgreater{} N};
  \item \textbf{horário:} \texttt{after @T}, \texttt{before @T}, \texttt{between @T @T}, \texttt{on [dias]};
  \item \textbf{solar:} \texttt{after sunset}, \texttt{before sunrise - 10m}, \texttt{between sunset sunrise} com deslocamento opcional;
  \item \textbf{gatilho:} \texttt{triggeredby "id"}.
\end{itemize}

Há suporte também para os operadores lógicos \texttt{not},
\texttt{and} e \texttt{or}, com precedência idêntica à linguagem
Python (\texttt{not} \textgreater{} \texttt{and} \textgreater{}
\texttt{or}) e parênteses para agrupar. A Listagem~\ref{lst:conditions} exemplifica algumas condições.

\begin{listing}[H]
    \centering
    \begin{minted}{swift}
if (
    light.sala == "on" and
    sensor.temperatura < 28 and
    between @06:00:00 @23:00:00 and
    not triggeredby "calor"
)
    \end{minted}
    \caption{Exemplos de condições.}
    \label{lst:conditions}
\end{listing}

A forma \texttt{between A B}, diferentemente de escrever
\texttt{after A and before B} como dois átomos, gera uma única
condição com ambos os limites, o que o Home Assistant interpreta como um
intervalo que pode atravessar dias distintos.

\subsection{Ações}

A seção \texttt{do} é uma sequência de comandos, executados em
ordem. As formas são chamadas de ação com \texttt{nome(arg, ...)}, esperas com \texttt{delay(DURAÇÃO)}, condicionais com \texttt{if}, \texttt{elif} e \texttt{else}, e blocos que continuam em caso de erro com \texttt{safe}. Na Listagem~\ref{lst:actions}, é possível observar algumas ações.

\begin{listing}[H]
    \centering
    \begin{minted}{swift}
do {
    turn_on(lampada, color=quente, transition=2s)   
    delay(5m)                                   
    if (triggeredby "manha") {
        turn_off(switch.luz)
    } else {
        toggle(switch.luz)
    }

    safe {
      say(tts.google, message="Aviso")
    }
}
    \end{minted}
    \caption{Exemplos de ações.}
    \label{lst:actions}
\end{listing}

As chamadas a serviços dependem diretamente do domínio do alvo. Isso significa que a mesma ação sobre dispositivos diferentes apresenta capacidades distintas. Por exemplo,
\texttt{turn\_on} sobre uma luz aceita \texttt{color}, enquanto que sobre um
interruptor, não. 

Por fim, como não é possível definir um vocabulário de ações e invocações diretas que cubra todos os casos, o Jass
permite invocar diretamente serviços com
\texttt{call("dominio.servico",\ ...)}, sendo a única ação isenta de validação de
argumentos, como demonstra a Listagem~\ref{lst:curated}.

\begin{listing}[H]
    \centering
    \begin{minted}{swift}
call("vacuum.send_command", target=vacuum.roomba, data={"command": "start"})
    \end{minted}
    \caption{Invocação direta de serviços.}
    \label{lst:curated}
\end{listing}

\subsection{Valores}

Os valores escritos em argumentos, metadados e nomes pertencem a um
conjunto fixo de tipos, resumido na Tabela~\ref{tab:value_types}. Tipos como
duração, horário, referência de entidade e
cor têm sintaxe própria.

\begin{table}[H]
\centering
\begin{tabular}{l l l}
\toprule
\textbf{Tipo} & \textbf{Exemplo} \\
\midrule
string & \texttt{"on"} \\
número & \texttt{30}, \texttt{0.09} \\
booleano & \texttt{True}, \texttt{False} \\
duração & \texttt{5m}, \texttt{1h30m}, \texttt{45s} \\
horário & \texttt{@07:00:00} \\
referência de entidade & \texttt{light.sala} \\
lista & \texttt{["off", "on"]} \\
dicionário & \texttt{\{"command": "start"\}} \\
cor & \texttt{rgb(255, 180, 80)} \\
\bottomrule
\end{tabular}
\caption{Tipos de valor da linguagem.}
\label{tab:value_types}
\end{table}

\subsection{Nomes e variáveis}

Para evitar repetição, o Jass
permite vincular nomes a valores com
\texttt{def NOME = VALOR}. Esses nomes são associações de tempo
de compilação, existindo apenas durante a compilação e sendo substituídos
pelo valor correspondente na saída. Cada uso é verificado por tipo exatamente como se o valor
estivesse escrito no lugar.

Há dois escopos: o de arquivo (global, visível a
qualquer automação declarada abaixo) e o de bloco (local ao
\texttt{\{ \}} em que o nome é declarado, com sombreamento, isto é,
um nome local oculta o de mesmo nome em escopo externo. A Listagem~\ref{lst:names} ilustra definições globais (\texttt{lampada},
\texttt{quente}) e local (\texttt{reforco}).

\begin{listing}[H]
    \centering
    \begin{minted}{swift}
def lampada = light.sala
def quente  = rgb(255, 180, 80)

automation "Luz aconchegante" {
    when (lampada changes to "on")
    do {
        turn_on(lampada, color=quente)
        def reforco = 100
        turn_on(lampada, brightness=reforco)
    }
}
    \end{minted}
    \caption{Exemplos de definições.}
    \label{lst:names}
\end{listing}

\section{Gramática Livre de Contexto (GLC)}

A sintaxe do Jass é definida no arqivo \href{https://github.com/jaimeadf-ufsm/elc408-compiladores/blob/main/trabalhos/t1/src/jass/grammar.lark}{\texttt{src/jass/grammar.lark}}, escrito na variante de EBNF
(\emph{Extended Backus--Naur Form}) adotada pelo Lark. Essa variante utiliza as seguintes notações para expressar a gramática:
\begin{itemize}
  \item as regras (não-terminais) são escritas em minúsculas e definidas
        com \texttt{:}, enquanto os terminais são escritos em maiúsculas;
  \item \texttt{\textbar{}} separa alternativas;
  \item \texttt{*} indica repetição (zero ou mais), \texttt{+} (uma ou mais) e
        \texttt{?} um elemento opcional (zero ou uma ocorrência);
  \item \texttt{(\ ...\ )} agrupa elementos;
  \item um terminal é definido por um literal entre aspas (\texttt{"when"}) ou
        por uma expressão regular entre barras (\texttt{/.../});
  \item \texttt{\%ignore} descarta um terminal do fluxo (espaços em branco e
        comentários) e \texttt{//} inicia um comentário no arquivo.
\end{itemize}

Além desses símbolos, existem duas diretivas que afetam apenas o formato da árvore sintática construída: o prefixo \texttt{?} em uma regra a embute,
removendo o seu nó da árvore quando ele possui um único filho, o que enxuga a
árvore resultante; e \texttt{->\ nome} atribui um rótulo à
alternativa, dando um nome estável ao nó por ela gerado. A Listagem~\ref{lst:grammar-rules} apresenta parcialmente a arquivo com a gramática, que pode ser visto na íntegra no repositório.

\begin{listing}[H]
    \centering
    \begin{minted}[fontsize=\scriptsize]{text}
start: (def_binding | automation)*

// Alias definitions
def_binding: DEF NAME ASSIGN value


// Automation definitions
automation: AUTOMATION STRING? LBRACE metadata* when_block? condition_block? do_block RBRACE
metadata: NAME ASSIGN value


// Triggers
when_block: WHEN LPAR trigger (OR trigger)* RPAR
trigger: trigger_expr (AS STRING)?
?trigger_expr: state_trigger
             | numeric_trigger
             | time_trigger
             | sun_trigger
state_trigger:   entity CHANGES (FROM value)? (TO value)? (FOR DURATION)?
numeric_trigger: entity (LT | GT) threshold (FOR DURATION)?
time_trigger:    TIME
sun_trigger:     sun_event


// Conditions
condition_block: IF LPAR bool_expr RPAR

...

// Values
TIME:     /@\d{2}:\d{2}:\d{2}/
DURATION: /\d+h(\d+m)?(\d+s)?|\d+m(\d+s)?|\d+s/
NUMBER:   /-?\d+\.\d+|-?\d+/
ENTITY:   /[a-z][a-z_]*\.[a-z0-9_]+/
NAME:     /[a-zA-Z_][a-zA-Z0-9_]*/
STRING:   /"(\\.|[^"\\])*"/


// Comments
COMMENT: /\/\/[^\n]*/
%ignore COMMENT
%ignore /[ \t\f\r\n]+/

    \end{minted}
    \caption{Conteúdo parcial da gramática com os não-terminais e os terminais.}
    \label{lst:grammar-rules}
\end{listing}

\section{Análise léxica}

A análise léxica converte o texto-fonte em uma sequência de
\emph{tokens}. No Jass, essa fase é inteiramente delegada ao
Lark, sendo as definições de terminais declaradas na própria
gramática, das quais o Lark deriva o \emph{lexer}. Comentários de linha (\texttt{//\ ...}) e espaços em branco são
descartados por diretivas \texttt{\%ignore}, não chegando ao
\emph{parser}. 

Quando o \emph{lexer} encontra um caractere que
não inicia nenhum terminal, o Lark lança a exceção
\texttt{UnexpectedCharacters}, que o Jass captura e
converte em um erro da categoria léxica, preservando a linha, a coluna e o caractere ofensor.
A Listagem~\ref{lst:lexical_error} mostra o
diagnóstico real para uma entrada que contém um \texttt{\$} inválido.

\begin{listing}[H]
    \centering
    \begin{minted}{swift}
// Caractere que não inicia nenhum token ('$').
automation "Caractere inválido" {
    when (binary_sensor.porta changes to "on")
    do {
        turn_on(light.sala$led)
    }
}
    \end{minted}
    \begin{minted}{text}
samples/errors/00-lexical.jass:5:27: lexical error: unexpected character '$'
          turn_on(light.sala$led)
                            ^
    \end{minted}
    \caption{Demonstração de erro léxico.}
    \label{lst:lexical_error}
\end{listing}







\section{Análise sintática}

A análise sintática verifica se a sequência de \emph{tokens} forma um
programa válido segundo a gramática e produz a árvore de
derivação. O Jass usa o \emph{parser} LALR(1) do Lark, construído
a partir da gramática. A análise opera em conjunto com o \textit{lexer} contextual,
que usa o estado do \emph{parser} para escolher, quando há ambiguidade
léxica, apenas entre os terminais válidos naquele ponto.


Quando o \emph{parser} recebe um \emph{token} que nenhuma entrada da
tabela aceita, o Lark sinaliza um \emph{token} inesperado, que o
Jass converte em um erro da categoria sintática, com posição e
a lista de \emph{tokens} esperados naquele ponto. Em vez de
parar no primeiro erro, o Jass aplica
recuperação em modo pânico, isto é, ao detectar a falha, registra o
erro, descarta \emph{tokens} até alcançar um ponto de
sincronização e retoma a análise a partir dele. Os pontos de
sincronização são as fronteiras de topo do programa, as
palavras-chave \texttt{automation} e \texttt{def}, que marcam o
início de uma nova unidade independente. Assim, um erro em uma automação
não contamina as seguintes, e uma única execução coleta todos
os erros. Alguns exemplos de erros sintáticos são apresentados na Listagem~\ref{lst:syntactic_errors}.

\begin{listing}[H]
    \centering
    \begin{minted}{swift}
// Falta a seção 'do' (obrigatória).
automation "Bloco do ausente" {
    when (binary_sensor.porta changes to "on")
    if (alarm_control_panel.alarmo == "disarmed")
}

// 'or' pendente: espera-se um gatilho após o último 'or'.
automation "Or pendente" {
    when (
        binary_sensor.sala changes to "on" as "sala" or
        binary_sensor.corredor changes to "on" as "corredor" or
    )
    do {
        turn_on(light.corredor)
    }
}

// Seções fora de ordem: 'if' antes de 'when'.
automation "Seções fora de ordem" {
    if (alarm_control_panel.alarmo == "disarmed")
    when (@08:00:00 as "manha")
    do { }
}
    \end{minted}
    \begin{minted}{text}
03-syntactic.jass:5:1: syntactic error: unexpected '}'; expected 'do'
  }
  ^
03-syntactic.jass:12:5: syntactic error: unexpected ')'; expected 'sunrise', 'sunset', a name, a time literal, ...
      )
      ^
03-syntactic.jass:20:5: syntactic error: unexpected 'if'; expected 'when', a name
      if (alarm_control_panel.alarmo == "disarmed")
      ^^
    \end{minted}
    \caption{Demonstração de erros sintáticos.}
    \label{lst:syntactic_errors}
\end{listing}

\section{Análise semântica}

A análise semântica é a fase em que o Jass concentra a maior
parte de suas verificações, percorrendo a AST de cada automação e coletando 
todos os erros. A análise é baseada em três pilares, tratados a seguir: a tabela de
símbolos, a verificação de tipos e a consistência ação e domínio. 
A Listagem~\ref{lst:semantic_errors} mostra alguns exemplos de erros semânticos.

\begin{listing}[H]
    \centering
    \begin{minted}{swift}
automation "Declaração" {
    when(light.sala changes)

    do {
        turn_on(luz)
    }
}

automation "Tipo" {
    when(light.sala changes)

    do {
        turn_on(light.sala, color=2s)
    }
}

automation "Ação e domínio" {
    when(light.sala changes)

    do {
        turn_on([light.sala, timer.sala])
    }
}         
    \end{minted}
    \begin{minted}{text}
05-semantic.jass:5:17: semantic error: undeclared name 'luz'
          turn_on(luz)
                  ^^^
05-semantic.jass:13:35: semantic error: argument 'color' expects a color, got a duration
          turn_on(light.sala, color=2s)
                                    ^^
05-semantic.jass:21:17: semantic error: all entities in a target list must share one domain (found: light, timer)
          turn_on([light.sala, timer.sala])
                  ^^^^^^^^^^^^^^^^^^^^^^
    \end{minted}
    \caption{Demonstração de erros semânticos.}
    \label{lst:semantic_errors}
\end{listing}


\subsection{Tabela de símbolos e escopos}

A tabela de símbolos registra os nomes vinculados por
\texttt{def}, associando a cada um o tipo inferido do valor e
o próprio valor. Os
tipos formam um conjunto fixo de categorias, como string, número,
booleano, duração, horário, entidade (que carrega também o domínio), lista, dicionário.

Os escopos são organizados em uma pilha aninhada que reflete a
estrutura do programa: o escopo de arquivo (global), o de cada
automação e o de cada bloco. A resolução de
um nome percorre a pilha do escopo atual para fora. Isso significa que uma definição local
de mesmo nome ``sombreia'' a externa. Nomes não resolvidos são erros. As referências dos
gatilhos declaradas com \texttt{as\ "id"} são coletadas por automação e
formam o conjunto contra o qual cada \texttt{triggeredby\ "id"} é
validado.

\subsection{Verificação de tipos}

Cada valor usado em um argumento, gatilho ou condição tem seu tipo
inferido e comparado ao tipo esperado naquele ponto. Quando o
valor é uma referência a uma definição, ela é resolvida ao seu
tipo antes de qualquer verificação.

\subsection{Consistência ação e domínio}

O significado de uma ação depende do domínio do alvo, isto é, do tipo do dispositivo indicado pelo primeiro parâmetro. O processo de análise ocorre de forma sequencial, começando pela determinação do domínio do alvo, com a exigência de que todas as entidades de um alvo em lista compartilhem o mesmo domínio. Em seguida, verifica-se se a ação está devidamente definida para ele. Por fim, validam-se os argumentos, garantindo que cada um seja aceito pela ação naquele domínio, apresente o tipo correto e respeite a obrigatoriedade, a ausência de duplicatas e a ligação posicional pela ordem. A única exceção a esse procedimento de análise é \texttt{call(...)}, que é isenta de validação.


\section{Geração de código}

A geração de código traduz a AST validada para a estrutura de saída do
Home Assistant. A AST é primeiro convertida
em um modelo de saída, um conjunto de dicionários e listas
Python, já no formato da plataforma. Em seguida é
serializado em YAML, com a biblioteca \texttt{ruamel.yaml}.

\section{Exemplos}

Os exemplos completos estão contidos no diretório \href{https://github.com/jaimeadf-ufsm/elc408-compiladores/tree/main/trabalhos/t1/samples}{\texttt{samples}}, com suas transpilações correspondentes. No entanto, para fins de demonstração no relatório, a Listagem~\ref{lst:example1_src} mostra a fonte de uma automação real, e a Listagem~\ref{lst:example1_yaml} exibe o YAML completo gerado para ela.

\begin{listing}[H]
    \centering
    \begin{minted}{swift}
def lampada   = light.sala
def amarelo   = rgb(255, 150, 60)
def tempo     = 5m

automation "Luz Sala" {
    description = "Ligar a luz quando ha movimento ou ao anoitecer"
    mode = "restart"

    when (
        binary_sensor.sala changes to "on" as "movimento" or
        sunset - 15m as "anoitecer"
    )

    if (
        between sunset sunrise and
        alarm_control_panel.casa == "disarmed"
    )

    do {
        turn_on(lampada, color=amarelo, transition=2s)

        if (triggeredby "movimento") {
            delay(tempo)
            turn_off(lampada)
        } else {
            say(tts.sala, player=media_player.sala, message="Bom anoitecer")
        }
    }
}
    \end{minted}
    \caption{Automação de exemplo em Jass.}
    \label{lst:example1_src}
\end{listing}

\begin{listing}[H]
    \centering
    \begin{minted}{yaml}
- id: "1780965949"
  alias: "Luz Sala"
  description: "Ligar a luz quando ha movimento ou ao anoitecer"
  mode: "restart"
  triggers:
    - trigger: state
      entity_id: binary_sensor.sala
      to: "on"
      id: "movimento"
    - trigger: sun
      event: sunset
      offset: "-00:15:00"
      id: "anoitecer"
  conditions:
    - condition: sun
      after: sunset
      before: sunrise
    - condition: state
      entity_id: alarm_control_panel.casa
      state: "disarmed"
  actions:
    - action: light.turn_on
      target:
        entity_id: light.sala
      data:
        rgb_color:
          - 255
          - 150
          - 60
        transition: 2
    - if:
        - condition: trigger
          id: "movimento"
      then:
        - delay: "00:05:00"
        - action: light.turn_off
          target:
            entity_id: light.sala
      else:
        - action: tts.speak
          target:
            entity_id: tts.sala
          data:
            media_player_entity_id: media_player.sala
            message: "Bom anoitecer"
    \end{minted}
    \caption{Saída completa em YAML para a Listagem~\ref{lst:example1_src}.}
    \label{lst:example1_yaml}
\end{listing}


\section{Conclusão}

O Jass cumpre os objetivos propostos, oferecendo uma sintaxe
declarativa e legível que esconde a complexidade do YAML de
automações do Home Assistant e permite a detecção de erros em tempo
de compilação. As falhas léxicas, sintáticas e semânticas são apresentadas com
categoria e localização, coletando vários problemas em uma única execução.
Todos os exemplos são transpilados para YAML válido e carregável.

\end{document}
